\chapter{Coagulation Factors}

\section*{What Is This Test?}
Coagulation factor assays measure the functional activity of individual coagulation factors in plasma. The coagulation cascade consists of a series of serine protease zymogens and cofactors that, upon activation, sequentially amplify the coagulation signal to produce a fibrin clot. The factors are designated by Roman numerals: factor I (fibrinogen), factor II (prothrombin), factor III (tissue factor), factor IV (calcium), factor V (labile factor), factor VII (stable factor), factor VIII (antihemophilic factor), factor IX (Christmas factor), factor X (Stuart-Prower factor), factor XI (plasma thromboplastin antecedent), factor XII (Hageman factor), and factor XIII (fibrin-stabilizing factor). Additionally, prekallikrein (Fletcher factor) and high-molecular-weight kininogen (HMWK, Fitzgerald factor) participate in the contact activation system. Individual factor assays are used to diagnose congenital factor deficiencies (hemophilia A = factor VIII, hemophilia B = factor IX, factor XI deficiency, and rare deficiencies of factors II, V, VII, X, XIII, and fibrinogen), to evaluate the severity of deficiency in the context of prolonged PT/aPTT, to monitor factor replacement therapy, and to detect acquired factor inhibitors (most commonly acquired hemophilia A due to anti-factor VIII autoantibodies). Factor assays are performed using clot-based, chromogenic, or immunologic methods.

\section*{Alternative Names}
Factor Assays; Clotting Factor Assays; Factor VIII Assay; Factor IX Assay; Factor Activity Tests; Coagulation Factor Levels; Individual Coagulation Factors

\section*{Principle}
Coagulation factor activity is most commonly measured by the one-stage clot-based assay, which is a modification of the PT (for factors II, V, VII, X) or aPTT (for factors VIII, IX, XI, XII, prekallikrein, HMWK). Patient plasma is serially diluted and mixed with factor-deficient plasma (commercially prepared plasma that is selectively deficient in a single factor but contains normal amounts of all other coagulation factors). The PT or aPTT is then measured. Because the deficient plasma provides all factors except the one being measured, the clotting time is proportional to the concentration of the patient's factor. The factor activity is determined by comparing the clotting time to a standard curve generated from serial dilutions of reference plasma (standardized against WHO international standards), and the result is expressed as a percentage of normal (normal = approximately 50--150\%) or in international units (IU/mL). Chromogenic substrate assays use a specific activator and a synthetic peptide substrate linked to a chromophore (typically para-nitroaniline, pNA). When the coagulation factor is activated, it cleaves the chromogenic substrate, releasing the chromophore, and the color intensity is measured spectrophotometrically. Chromogenic assays are more precise for low factor levels and are used for factor VIII in patients with mild hemophilia A (to distinguish hemophilia A from the cross-reacting material variant) and for monitoring factor replacement therapy. Factor XIII is measured separately using a clot solubility test (qualitative) or a quantitative chromogenic assay. Fibrinogen (factor I) is measured by the Clauss method. Immunologic (ELISA) methods measure factor antigen (protein concentration) and are used when functional activity and antigen levels diverge (cross-reacting material-positive [CRM+] hemophilia: normal antigen, reduced activity).

\section*{History}
The coagulation cascade was progressively elucidated through the 20th century. Factor VIII deficiency was recognized as the cause of hemophilia in the 19th century (the ``royal disease'' in Queen Victoria's descendants). Factor IX was identified in 1952 when a patient named Stephen Christmas was found to have a factor distinct from factor VIII, leading to the designation ``Christmas disease'' (hemophilia B). The remaining coagulation factors and the cascade model were described by Davie, Ratnoff, and Macfarlane in 1964 (the ``waterfall'' or ``cascade'' hypothesis). The development of specific factor-deficient plasmas and lyophilized reference standards in the 1960s--1970s enabled the one-stage clot-based assay that remains the standard today. The introduction of recombinant factor VIII (1993) and factor IX (1997) revolutionized hemophilia treatment. Chromogenic assays, genetic testing, and extended half-life factor concentrates represent the modern era of factor assay development.

\section*{Why This Test Is Ordered}
\begin{itemize}
  \item Diagnosis of congenital factor deficiencies --- hemophilia A (factor VIII deficiency), hemophilia B (factor IX deficiency), factor XI deficiency, and rare deficiencies of factors II, V, VII, X, XIII, and fibrinogen
  \item Investigation of prolonged PT (isolated prolonged PT suggests factor VII deficiency) or prolonged aPTT (isolated prolonged aPTT suggests factor VIII, IX, XI, or XII deficiency)
  \item Determination of factor deficiency severity --- hemophilia A/B severity: severe <1\% factor activity, moderate 1--5\%, mild 5--40\%
  \item Preoperative evaluation in patients with personal or family history of bleeding disorders
  \item Diagnosis of von Willebrand disease --- factor VIII activity is reduced due to decreased vWF-mediated factor VIII stabilization
  \item Detection of acquired coagulation factor inhibitors --- acquired hemophilia A (anti-factor VIII autoantibody), which presents with severe bleeding in adults without prior bleeding history
  \item Monitoring of factor replacement therapy (factor VIII or IX concentrates, recombinant factors, cryoprecipitate, fresh frozen plasma) in hemophilia patients
  \item Evaluation of vitamin K status --- factors II, VII, IX, X (and proteins C, S) require vitamin K for gamma-carboxylation
  \item Assessment of liver synthetic function --- factors II, V, VII, IX, X are synthesized in hepatocytes; factor V levels help distinguish liver disease (low factor V) from vitamin K deficiency or warfarin effect (normal factor V)
  \item Determination of factor levels before major surgery in liver disease
  \item Investigation of unexplained, recurrent venous thromboembolism with prolonged PT or aPTT that corrects on mixing study
\end{itemize}

\section*{Is This a Primary or Secondary Test}
Coagulation factor assays are secondary, specific tests ordered when screening PT and aPTT are prolonged and a mixing study demonstrates correction (factor deficiency pattern).

\section*{How This Test Is Performed}
A venous blood sample is collected in a light blue-top tube containing 3.2\% sodium citrate (9:1 ratio), filled exactly to the line. The sample is centrifuged at 1,500--2,500 $\times$ g for 15 minutes to obtain platelet-poor plasma. For the one-stage assay, patient plasma is serially diluted, mixed with factor-deficient plasma specific to the factor being measured, and the PT (for extrinsic pathway factors) or aPTT (for intrinsic pathway factors) is measured on an automated coagulation analyzer. The clotting time is compared to a standard curve. Factor XIII is measured separately by clot solubility in 5M urea or 1\% monochloroacetic acid (screening test --- a normal clot remains intact in urea) or by quantitative chromogenic assay. Factor assays are typically batched and may take 2--4 hours to complete. Venipuncture takes less than 5 minutes.

\section*{What Preparation Is Needed}
No special preparation is required. Fasting is not necessary. The patient must inform the healthcare provider of all medications, especially: anticoagulants (warfarin, heparin, DOACs -- all lower vitamin K-dependent factor levels or falsely affect clot-based assays), factor replacement products (factor VIII/IX concentrates, desmopressin/DDAVP which releases factor VIII and vWF), recent blood transfusion, fresh frozen plasma, cryoprecipitate, or prothrombin complex concentrate administration. Pregnancy (increases factor VIII), acute illness (increases factor VIII as an acute-phase reactant), and liver disease should be disclosed. Testing should ideally be performed when the patient is not on anticoagulation; if impossible, the interpreting physician must account for anticoagulant effects on each factor.

\section*{Contraindications and Precautions}
No absolute contraindications. Critical pre-analytical requirements: the light blue-top tube must be filled exactly to the line; hemolyzed, clotted, or lipemic specimens are unacceptable; traumatic venipuncture can activate coagulation and affect factor levels; heparin contamination from lines causes falsely low factor activity in clot-based assays; warfarin and DOACs interfere with clot-based factor assays -- factor VIII, IX, XI, XII can be reliably measured on warfarin (they are not vitamin K-dependent), but factors II, VII, IX, X are affected; factor VIII is an acute-phase reactant and is normally elevated in acute illness, pregnancy, inflammation, and stress; elevated factor VIII can mask mild von Willebrand disease or hemophilia A carrier state; factor assays require comparison to a reference plasma standard curve, which must be traceable to WHO international standards; quality control materials at normal and abnormal (low) levels must be run with each batch; lithium-heparin plasma, serum, or serum separator tubes are NOT acceptable (citrate is the only acceptable anticoagulant); plasma should be separated and frozen within 4 hours if testing is not immediate; frozen plasma is stable for 2--4 weeks at -20\textdegree{}C; many factors are labile and degrade with repeated freeze-thaw cycles, especially factor V and factor VIII. Factor VIII and vWF assays should be performed on fresh plasma whenever possible because cold storage can activate or precipitate these proteins.

\section*{Risks}
Minimal risks of venipuncture: mild pain, bruising, hematoma, lightheadedness, and rarely infection or nerve injury. Patients with bleeding disorders or on anticoagulants have increased risk of hematoma; firm pressure should be applied for 5 minutes after venipuncture.

\section*{What Happens After the Test}
Results are typically available within 2--4 hours. For factor deficiency diagnosis, the specific factor activity level determines the diagnosis and severity: severe (<1\%), moderate (1--5\%), mild (5--40\%) for hemophilia. Results are interpreted with PT, aPTT, and mixing study. For acquired factor inhibitor evaluation, a Bethesda assay quantifies the inhibitor titer and is performed separately. For liver disease, a factor V level (low) distinguishes from vitamin K deficiency or warfarin effect (factor V is normal because it is not vitamin K-dependent). For hemophilia monitoring, factor VIII or IX levels are checked after factor concentrate infusion to verify adequate dosing. Perioperative management uses factor levels to ensure hemostasis (target: at least 50--80\% for major surgery, 30--50\% for minor procedures).

\section*{Understanding Results}

\subsection*{Below Normal (Factor Deficiency)}
\begin{itemize}
  \item \textbf{Hemophilia A (factor VIII deficiency):} X-linked recessive. Factor VIII activity <1\% (severe -- spontaneous hemarthroses, intramuscular bleeding, intracranial hemorrhage), 1--5\% (moderate), 5--40\% (mild). Prolonged aPTT, normal PT, normal platelet count. Diagnosed by factor VIII assay. Approximately 30\% of severe hemophilia A patients develop factor VIII inhibitors (alloantibodies), a serious complication detected by Bethesda assay.
  \item \textbf{Hemophilia B (factor IX deficiency):} X-linked recessive. Identical clinical presentation to hemophilia A. Diagnosed by factor IX assay. Factor IX inhibitor development is less common than factor VIII inhibitor.
  \item \textbf{Hemophilia C (factor XI deficiency):} Autosomal recessive (more common in Ashkenazi Jewish and Iraqi Jewish populations, prevalence up to 8\%). Bleeding is milder than hemophilia A/B, typically provoked by surgery or trauma (especially involving sites with high fibrinolytic activity: oral cavity, urogenital tract). Spontaneous bleeding is uncommon. Prolonged aPTT, normal PT. Factor XI level correlates poorly with bleeding tendency.
  \item \textbf{Factor VII deficiency (extrinsic pathway):} Autosomal recessive. Isolated prolonged PT with normal aPTT. Most common rare bleeding disorder. Bleeding severity varies; severe deficiency (<1\%) causes CNS and GI bleeding in infancy; mild deficiency (5--40\%) may have minimal bleeding. Some patients with factor VII activity <1\% have no bleeding.
  \item \textbf{Factor X deficiency (common pathway):} Autosomal recessive. Prolonged PT and aPTT. Severe deficiency causes hemarthroses, mucocutaneous bleeding, and intracranial hemorrhage. Acquired factor X deficiency occurs in AL amyloidosis (factor X binds to amyloid fibrils).
  \item \textbf{Factor V deficiency (common pathway):} Autosomal recessive. Prolonged PT and aPTT. Mucocutaneous bleeding, epistaxis, menorrhagia. Combined factor V and factor VIII deficiency is a distinct autosomal recessive disorder (defect in LMAN1 or MCFD2 genes). Factor V is not vitamin K-dependent; levels distinguish liver disease (low factor V) from vitamin K deficiency/warfarin (normal factor V).
  \item \textbf{Factor II (prothrombin) deficiency:} Autosomal recessive. Prolonged PT and aPTT. Severe deficiency (<1\%) causes severe bleeding in infancy; milder deficiency causes post-traumatic/surgical bleeding.
  \item \textbf{Factor XIII deficiency:} Autosomal recessive. PT and aPTT are normal (clot forms but is unstable because cross-linking is absent). Presents with delayed bleeding (hours to days after injury), umbilical stump bleeding (characteristic in neonates), intracranial hemorrhage, and recurrent miscarriage. Diagnosed by clot solubility test (abnormal clot dissolves in 5M urea) or quantitative assay. Normal PT/aPTT with clinical bleeding should raise suspicion.
  \item \textbf{Factor XII, prekallikrein, or HMWK deficiency:} Prolonged aPTT but NO clinical bleeding. Important to recognize to avoid unnecessary investigation and treatment. Factor XII deficiency may be paradoxically associated with thrombosis (limited evidence).
  \item \textbf{Vitamin K deficiency:} Decreased activity of factors II, VII, IX, X and proteins C and S. Prolonged PT (factor VII has the shortest half-life, 4--6 hours), then prolonged aPTT as other factors decline. Causes: malnutrition, prolonged antibiotics, biliary obstruction, malabsorption, warfarin therapy.
  \item \textbf{Warfarin therapy:} Decreased activity of vitamin K-dependent factors (II, VII, IX, X, proteins C, S). Prolonged PT/INR, variably prolonged aPTT. Factor VII decreases first, followed by protein C (also short half-life, causing transient hypercoagulable state if not bridged with heparin). Factor V is normal.
  \item \textbf{Liver disease:} Decreased synthesis of all coagulation factors except factor VIII (synthesized by sinusoidal endothelial cells) and von Willebrand factor. Factor V is decreased (distinguishes from vitamin K deficiency/warfarin). Prolonged PT and aPTT, with low factor V confirming hepatic synthetic failure.
  \item \textbf{DIC:} Consumptive decrease in all coagulation factors (especially factors V, VIII, and fibrinogen). Prolonged PT and aPTT, thrombocytopenia, elevated D-dimer, low fibrinogen.
  \item \textbf{Acquired coagulation factor inhibitors:} Autoantibodies against specific factors, most commonly factor VIII (acquired hemophilia A). Presents with severe, life-threatening bleeding in adults without prior bleeding history; prolonged aPTT, factor VIII inhibitor by Bethesda assay. Also anti-factor V, anti-factor II.
  \item \textbf{von Willebrand disease:} Decreased factor VIII activity (vWF stabilizes factor VIII). Diagnosed by vWF antigen, ristocetin cofactor activity, and vWF multimer analysis.
\end{itemize}

\subsection*{Above Normal (Elevated Factor Activity)}
\begin{itemize}
  \item \textbf{Factor VIII elevation:} Most common coagulation factor elevation. Acute-phase reactant elevated in infection, inflammation, pregnancy, estrogen therapy, malignancy, surgery, trauma, stress, and hyperthyroidism. Associated with increased risk of venous thromboembolism. Elevated factor VIII can mask mild hemophilia A or VWD on aPTT screening.
  \item \textbf{von Willebrand factor elevation:} Similar to factor VIII; acute-phase reactant.
  \item \textbf{Other factor elevations:} Factor VII, fibrinogen, and vWF are all acute-phase reactants and are commonly elevated in inflammatory states. Factor IX and factor II elevations are uncommon but may contribute to thrombosis.
  \item \textbf{Factor V Leiden:} A genetic mutation (factor V Leiden, R506Q) producing activated protein C resistance; factor V activity is normal, but the protein is resistant to inactivation by protein C. NOT detected by standard factor V activity assay; diagnosed by activated protein C resistance assay or DNA testing.
\end{itemize}

\section*{Normal Results}
\begin{itemize}
  \item \textbf{All coagulation factors:} 50--150\% activity (0.50--1.50 IU/mL) in healthy adults
  \item \textbf{Vitamin K-dependent factors (II, VII, IX, X):} 50--140\%
  \item \textbf{Factor VIII and vWF:} Wide normal range (50--200\%) due to acute-phase reactivity; blood type influences levels --- individuals with blood group O have 25--30\% lower vWF and factor VIII than non-O blood groups
  \item \textbf{Newborns:} Vitamin K-dependent factors (II, VII, IX, X) and factors XI, XII, prekallikrein, HMWK are physiologically low at birth (20--50\%), normalizing by 3--6 months. Factor VIII, V, XIII, and fibrinogen are at adult levels at birth.
  \item \textbf{Pregnancy (third trimester):} Factor VIII, vWF, factor VII, and fibrinogen are elevated. Factors II, X slightly elevated. Protein S is decreased (60--70\% of normal).
  \item \textbf{Elderly:} Factor VIII, vWF, factor VII, and fibrinogen increase with age.
  \item \textbf{Blood group O:} vWF and factor VIII are 25--30\% lower than other blood groups.
\end{itemize}

\section*{Follow-up Tests}
\begin{itemize}
  \item \textbf{Bethesda assay for factor inhibitors:} Quantifies factor inhibitor titer (1 Bethesda unit = amount of inhibitor that neutralizes 50\% of factor activity when patient plasma is incubated with normal plasma for 2 hours at 37\textdegree{}C)
  \item \textbf{Genetic testing:} For specific mutations in factor genes (F8 gene for hemophilia A, F9 for hemophilia B, F11 for factor XI, etc.)
  \item \textbf{von Willebrand factor panel:} vWF antigen, vWF ristocetin cofactor activity, factor VIII activity, vWF multimers, and ristocetin-induced platelet aggregation (RIPA) for VWD diagnosis
  \item \textbf{Chromogenic factor VIII assay:} To distinguish hemophilia A from CRM+ variant (normal clotting assay but reduced chromogenic activity)
  \item \textbf{Mixing studies for PT and aPTT:} To confirm factor deficiency pattern before specific factor assays
  \item \textbf{Liver function tests, albumin, INR:} For assessment of hepatic synthetic function
  \item \textbf{Serum vitamin K level or PIVKA-II (proteins induced by vitamin K absence):} When vitamin K deficiency is suspected
  \item \textbf{Thromboelastography (TEG) or rotational thromboelastometry (ROTEM):} Global assessment of hemostasis including factor function
\end{itemize}

\section*{References}
\begin{enumerate}
  \item Tierney LM, Saint S, Drazen JM. CURRENT Medical Diagnosis \& Treatment. McGraw-Hill; 2023.
  \item Wallach J. Interpretation of Diagnostic Tests. 11th ed. Wolters Kluwer; 2022.
  \item Fischbach FT, Dunning MB. A Manual of Laboratory and Diagnostic Tests. 11th ed. Wolters Kluwer; 2023.
  \item Burtis CA, Ashwood ER, Bruns DE. Tietz Textbook of Clinical Chemistry and Molecular Diagnostics. 7th ed. Elsevier; 2023.
  \item Kaushansky K, Lichtman MA, Prchal JT, et al. Williams Hematology. 10th ed. McGraw-Hill; 2021.
  \item Kitchen S, McCraw A, Echenagucia M. Diagnosis of hemophilia and other bleeding disorders. A laboratory manual. 2nd ed. World Federation of Hemophilia; 2010.
  \item Tripodi A, Chantarangkul V, Mannucci PM. Acquired coagulation disorders: clinical and laboratory aspects. Hematology Am Soc Hematol Educ Program. 2006:99--105.
  \item Di Minno MND, Navone R, Di Minno G. Measurement of coagulation factor levels. In: Key NS, ed. Hemophilia and von Willebrand Disease. Elsevier; 2018.
\end{enumerate}

\clearpage
